\documentclass[11pt,oneside]{article}
\usepackage[margin=1in]{geometry}
\usepackage{hyperref}
\usepackage{graphicx}
\usepackage{xcolor}
\usepackage{fancyhdr}
\usepackage{booktabs}
\usepackage{array}
\usepackage{listings}
\usepackage{xparse}

% Code formatting
\lstset{
  basicstyle=\ttfamily\small,
  breaklines=true,
  frame=single,
  backgroundcolor=\color{gray!10},
  keywordstyle=\color{blue},
  commentstyle=\color{gray},
  stringstyle=\color{red},
  showstringspaces=false,
  language=json
}

% Header and Footer
\pagestyle{fancy}
\fancyhf{}
\fancyhead[L]{\textbf{AEGIS PRD v1.0.0}}
\fancyhead[R]{\today}
\fancyfoot[C]{\thepage}

% Title customization
\title{\textbf{\Large MASTER PRODUCT REQUIREMENTS DOCUMENT (PRD)}\\[0.3cm]{\normalsize Project: AEGIS (Working Title)}}
\author{Lead Architect: Surya}
\date{Document Version: 1.0.0 (Master Blueprint)}

\begin{document}

\maketitle

\section*{Platform}
Cross-platform Mobile (iOS/Android) \& Web Application

\tableofcontents
\newpage

%---------------------------------------------------------------------------
\section{Executive Summary}
%---------------------------------------------------------------------------

Aegis is an AI-driven ``Operating System for Life.'' It resolves app fragmentation by merging task management, physical fitness, nutrition, financial bookkeeping, deep-work enforcement, and social gamification into a single, unified ecosystem. The platform is orchestrated by a self-hosted, hyper-personalized Large Language Model (LLM) that acts as a voice-activated router, converting natural language directly into structured database logs while adopting dynamic personas to motivate the user.

%---------------------------------------------------------------------------
\section{Product Vision \& Problem Statement}
%---------------------------------------------------------------------------

\subsection*{The Problem}
High-performing individuals suffer from the ``Context Switching Tax.'' Managing complex coding sprints, university coursework, physical health, family logistics, and personal finances requires jumping between disconnected silos. These silos prevent a holistic understanding of daily productivity.

\subsection*{The Solution}
A unified MongoDB-backed architecture where a single spoken sentence updates the user's financial ledger, task list, and physical health dashboard simultaneously. All actions feed into a centralized ``Productivity Score'' and live social leaderboard to drive retention through gamification and extreme accountability.

%---------------------------------------------------------------------------
\section{Target User Personas}
%---------------------------------------------------------------------------

Aegis targets users who experience high friction managing complex, multi-faceted lives:

\begin{itemize}
  \item \textbf{The Deep-Work Developer/Student:} Needs to manage intense academic sprints, track WakaTime coding hours, and utilize hardcore app blockers during hackathon prep.
  
  \item \textbf{The Co-op Coordinator:} Requires a shared environment for household sync, managing joint grocery budgets, and coordinating family logistics (e.g., managing morning school routines in Thane).
  
  \item \textbf{The Optimizer:} Wants frictionless entry. Relies on Vision AI to snap photos of meals for instant macro-tracking and voice commands for instant expense logging.
  
  \item \textbf{The External Accountability Seeker:} Users who actively want an AI to adopt a harsh persona to roast them when screen time spikes, acting as an uncompromising digital coach.
\end{itemize}

%---------------------------------------------------------------------------
\section{Core Feature Specifications}
%---------------------------------------------------------------------------

\subsection{The ``Chameleon'' AI Agent (The Brain)}

The primary UI is the AI itself, accessible instantly.

\begin{itemize}
  \item \textbf{Instant Access:} Triggered via native iOS/Android home-screen widgets or physical Action Button mapping. The listening UI must load in $< 500$ms.
  
  \item \textbf{Dynamic Personas:} The AI injects user preferences dynamically. It can act as a Gen Z hype-man, a formal assistant, or adopt a fictional persona (e.g., Kratos from God of War delivering tough love regarding broken C pointers or skipped workouts).
  
  \item \textbf{Multi-Intent Routing:} Parses complex commands (``Log 500 rupees for coffee, I ran 2km, and add a Linear Algebra review to my sprint'') into structured JSON payloads.
\end{itemize}

\subsection{Unified Task \& Sprint Engine}

\begin{itemize}
  \item \textbf{Project Sprints:} Converts a massive goal (like a Yugastr hackathon submission) into daily, chronological micro-tasks.
  
  \item \textbf{Household Sync:} Shared database collections for group lists and delegated chores.
  
  \item \textbf{Automated Rollover:} Uncompleted tasks roll over to the next day, accompanied by an AI-generated motivational nudge or roast.
\end{itemize}

\subsection{Holistic Health \& Fitness}

\begin{itemize}
  \item \textbf{Vision-to-Log:} Users upload a meal photo; the Gemini Vision API returns estimated calories and macronutrients.
  
  \item \textbf{Fitness Ledger:} Granular database logging for workouts, tracking sets, reps, weight, and cardio distance.
\end{itemize}

\subsection{The ``Doom Protocol'' (Screen Time \& Deep Work)}

\begin{itemize}
  \item \textbf{App Blocking:} Integrates with OS-level APIs (Apple Family Controls / Android UsageStats) to hard-block distracting apps.
  
  \item \textbf{IDE Integrations:} Webhooks connected to WakaTime/GitHub to automatically log time spent actively coding.
  
  \item \textbf{Cognitive Barriers:} High-stakes mode where failing a priority task locks distracting apps behind a cognitive barrier (e.g., forcing the user to calculate a matrix row reduction to unlock Instagram).
\end{itemize}

\subsection{Financial Ledger \& Social Arena}

\begin{itemize}
  \item \textbf{Voice Ledger:} Natural language processing of expenses instantly categorizes and deducts amounts from a dynamic budget.
  
  \item \textbf{The RPG Score:} A proprietary algorithm weighting all actions into a single daily Productivity Score.
  
  \item \textbf{Live Sockets:} A real-time activity stream broadcasting friend milestones and challenge updates instantly.
\end{itemize}

%---------------------------------------------------------------------------
\section{System Architecture \& Tech Stack}
%---------------------------------------------------------------------------

\begin{table}[h]
\centering
\begin{tabular}{|p{2.5cm}|p{3cm}|p{5cm}|}
\hline
\textbf{Layer} & \textbf{Technology} & \textbf{Purpose} \\
\hline
Monorepo Strategy & Turborepo & Shares types, UI components, and API logic across Web and Mobile. \\
\hline
Web Frontend & Next.js (App Router), Tailwind, Zustand & SEO-friendly dashboards and complex sprint planning interfaces. \\
\hline
Mobile Client & React Native, Expo, NativeWind & Cross-platform iOS/Android deployment with native OS integrations. \\
\hline
Backend API & Next.js API Routes & Serverless execution of standard CRUD operations and database reads. \\
\hline
Real-Time Engine & Node.js / Socket.io & Manages live leaderboards, activity feeds, and instant challenge updates. \\
\hline
Primary Database & MongoDB & Document-based structure to handle polymorphic action logs dynamically. \\
\hline
AI Infrastructure & Llama 3 (8B) / Ollama / Groq & Self-hosted model for secure, zero-latency text-to-JSON routing. \\
\hline
Vision API & Gemini Flash / Claude & Offloaded processing for complex image-to-text tasks (food/receipt scanning). \\
\hline
\end{tabular}
\caption{Core Technology Stack for AEGIS Platform}
\end{table}

%---------------------------------------------------------------------------
\section{Database Schema (MongoDB)}
%---------------------------------------------------------------------------

The system relies on a polymorphic \texttt{ActionLog} collection to maintain a unified timeline.

\subsection*{Users Collection}

\begin{lstlisting}
{
  "_id": "ObjectId",
  "name": "Surya",
  "productivityScore": 1450,
  "personaSettings": {
    "character": "Kratos",
    "tone": "Tough Love",
    "roastLevel": 9
  },
  "integrations": {
    "wakatimeKey": "String",
    "githubToken": "String"
  },
  "friends": ["ObjectId"],
  "householdId": "ObjectId"
}
\end{lstlisting}

\subsection*{ActionLogs Collection (The Master Ledger)}

\begin{lstlisting}
{
  "_id": "ObjectId",
  "userId": "ObjectId",
  "timestamp": "ISODate",
  "domain": "String",
  "payload": {
     // Dynamic Object based on Domain.
     // Example for TASK:
     // {"title": "Data Structures Lab", "status": "completed"}
     // Example for EXPENSE:
     // {"amount": 500, "category": "Food"}
  }
}
\end{lstlisting}

\noindent
\textbf{Note:} Domain enumeration: \texttt{TASK}, \texttt{EXPENSE}, \texttt{FITNESS}, \texttt{DEEP\_WORK}

%---------------------------------------------------------------------------
\section{Core User Flow: The Voice-to-Database Pipeline}
%---------------------------------------------------------------------------

\begin{enumerate}
  \item \textbf{Trigger:} User taps the mobile widget.
  
  \item \textbf{Capture:} React Native voice module records the command.
  
  \item \textbf{Transcription:} Fast cloud API or on-device model converts audio to text.
  
  \item \textbf{Inference:} Text is routed to the Llama 3 server, injected with the user's \texttt{personaSettings}.
  
  \item \textbf{Routing:} The LLM returns a structured JSON array of actions and custom dialogue.
  
  \item \textbf{Database Write:} Next.js backend parses the JSON and inserts documents into the \texttt{ActionLogs} collection.
  
  \item \textbf{Real-Time Broadcast:} The backend recalculates the RPG Productivity Score and fires a Socket.io event.
  
  \item \textbf{Client Update:} The mobile UI instantly renders the AI's dialogue, updates personal stats, and pushes the event to friends' live feeds.
\end{enumerate}

%---------------------------------------------------------------------------
\section{Development Milestones \& Go-To-Market}
%---------------------------------------------------------------------------

\subsection{Phase 1: The Core Loop (Weeks 1--4)}

\begin{itemize}
  \item Initialize Turborepo with Next.js and Expo.
  \item Set up MongoDB schemas and establish the Groq/Llama 3 routing logic for Task and Expense generation.
\end{itemize}

\subsection{Phase 2: Deep Work \& Social Engine (Weeks 5--8)}

\begin{itemize}
  \item Implement WakaTime/GitHub webhooks.
  \item Build the unified Productivity Score algorithm and integrate Socket.io for the live activity feed.
\end{itemize}

\subsection{Phase 3: Native Integrations (Weeks 9--12)}

\begin{itemize}
  \item Integrate Gemini for Vision AI (food logging).
  \item Implement custom Expo modules for Screen Time blocking (Family Controls/UsageStats).
\end{itemize}

\subsection{Go-To-Market Strategy: The Trojan Horse}

Launch V1.0 strictly targeting the developer/student niche as a ``Deep Work \& Hackathon Planner.'' Highlight the ruthless screen-time blocking and AI roasting. Once the core loop is proven and retention is high, expand into the Healthify and Co-op features for broader market adoption.

%---------------------------------------------------------------------------
\section*{Conclusion}
%---------------------------------------------------------------------------

The AEGIS platform represents a comprehensive reimagining of personal productivity through integrated AI orchestration. The modular tech stack and phased rollout strategy position the product for rapid iteration, market validation, and sustainable scaling.

\end{document}
